% 同策略蒸馏：原理与两次训练。中文版（译自 opd-report-en.tex）。
% Build: cd .report && tectonic -X compile opd-report-zh.tex
% 版式：warmreport.sty。中文字体用 macOS 自带的宋体、楷体和冬青黑体。
\documentclass[11pt,a4paper,fontset=none]{ctexart}
\usepackage{warmreport}
\setCJKmainfont{Songti SC}[BoldFont={Songti SC Bold},ItalicFont={Kaiti SC}]
\setCJKsansfont{Hiragino Sans GB W3}[BoldFont={Hiragino Sans GB W6}]
\setCJKmonofont{Hiragino Sans GB W3}
\renewcommand{\wrshorttitle}{从第一性原理理解同策略蒸馏}
\hypersetup{pdftitle={从第一性原理理解同策略蒸馏}}
\renewcommand{\algoname}{算法}
\renewcommand{\algorithmicrequire}{\textbf{输入：}}

\newcommand{\pis}{\pi_\theta}
\newcommand{\pit}{\pi_T}
\newcommand{\KL}{\mathrm{KL}}
\newcommand{\E}{\mathbb{E}}

% 图中标签（curves.tex）。
\newcommand{\LblKL}{教师 KL（nats/token）}
\newcommand{\LblReward}{奖励（任务解决率）}
\newcommand{\LblLength}{回合长度（token）}
\newcommand{\LblEntropy}{学生熵（nats/token）}
\newcommand{\LblStep}{训练步数，阶段 1--4 首尾相接}
\newcommand{\LblQwenThree}{Qwen3-8B $\to$ Qwen3-1.7B}
\newcommand{\LblQwenThreeFive}{Qwen3.5-9B $\to$ Qwen3.5-2B}

\begin{document}

\wrtitle{技术报告 · hands-on-policy-distillation}
  {从第一性原理理解同策略蒸馏}
  {两个小型 Qwen 模型学会使用工具}
  {2026 年 10 月 · 一块 A100-80GB 上的 Qwen3-8B $\to$ Qwen3-1.7B 与
  Qwen3.5-9B $\to$ Qwen3.5-2B}

\begin{wrabstract}{摘要}
我们在一块 A100-80GB GPU 上，用同一套代码和超参数，把 Qwen3-8B 蒸馏进
Qwen3-1.7B，把 Qwen3.5-9B 蒸馏进 Qwen3.5-2B，共四个阶段：小学数学应用题、
计算器工具、多轮回合中的三种工具，以及蒸馏加任务奖励。我们从“学生该模仿什么、
在哪里模仿”出发，推导出同策略蒸馏就是一种策略梯度：每个 token 的奖励等于教师
对学生所写 token 的对数概率减去学生自己的对数概率。每个蒸馏阶段的教师 KL 都在
下降；在完全不用奖励的情况下，任务成功率也随之上升：计算器任务上从 0.47 升到
1.00、从 0.65 升到 0.90，三工具任务上从 0.66 升到 0.75、从 0.75 升到 0.96。
再通过 GRPO 加入任务奖励，没有带来可测量的变化，因为分别有 98.5\% 和 89\% 的
采样组得到完全相同的奖励。Qwen3.5 的训练用了 2.5 倍的 GPU 时间，原因是它的
回合长 2--5 倍。
\end{wrabstract}

\begin{tcbraster}[raster columns=4, raster equal height,
    raster column skip=3mm]
\begin{stat}{wr-terracotta}{0.47 $\to$ 1.00}
Qwen3 学生的计算器任务解决率，未用奖励
\end{stat}
\begin{stat}{wr-ochre}{0.75 $\to$ 0.96}
Qwen3.5 学生的三工具任务解决率，未用奖励
\end{stat}
\begin{stat}{wr-wine}{98.5\%}
Qwen3 的 GRPO 组奖励完全相同：没有信号
\end{stat}
\begin{stat}{wr-sage}{2.5$\times$}
Qwen3.5 的 GPU 时间，因其回合长 2--5 倍
\end{stat}
\end{tcbraster}

\section{引言}

能调用工具的小模型部署成本低，但它得学会：什么时候调用哪个工具、调用怎么写、
拿到结果之后怎么办。教会它有三条路：在更大的教师模型写出的回合上做监督微调
（异策略蒸馏；\citealp{kim2016seqkd}）；用任务奖励做强化学习，例如 GRPO
\citep{shao2024deepseekmath}；以及同策略蒸馏（on-policy distillation）——
学生自己写回合，教师给每个 token 打分
\citep{agarwal2024gkd,gu2024minillm,lu2025onpolicy}。Qwen3 的小模型本身就是
用同策略蒸馏从更大的模型后训练出来的 \citep{yang2025qwen3}。

本报告先从第一性原理解释第三条路为什么可行（第 \ref{sec:principle} 节），
再说明我们如何训练（第 \ref{sec:setup} 节），以及两次训练的结果
（第 \ref{sec:results} 节）。两次训练用的是同一条流水线（在 TRL 的 GRPO
训练器的优势中加入教师项）和同一套配置，只有模型不同。

\section{从第一性原理看核心思想}
\label{sec:principle}

\paragraph{设定}
语言模型是定义在 token 上的策略。状态 $s_t$ 是学生写第 $t$ 个 token 之前
上下文中的全部内容：提示（系统提示、工具说明、任务）、学生先前写下的 token，
以及它调用过的工具返回的结果。动作是下一个 token $y_t$。学生 $\pis$ 参与训练，
教师 $\pit$ 保持冻结。工具结果由环境写入，而不是学生选择的，所以它属于状态，
永远不是动作。

我们希望学生在每个状态下都像教师那样选择 token。衡量两个下一 token 分布之间
差异的自然选择是 KL 散度，于是只剩两个问题：在哪些状态上衡量，取哪个方向。

\paragraph{在哪些状态上：学生自己的状态}
在教师写的回合上微调，训练的是教师会访问的状态。可到了测试时，生成的是学生，
访问的是它自己的状态。只要写出一个教师不会写的 token，学生就进入了训练样本
从未覆盖的状态，在那里更容易再次出错：误差会累积。在模仿学习中，若训练分布上
每步的错误率为 $\epsilon$，$T$ 步的总代价可能按 $O(\epsilon T^2)$ 增长
\citep{ross2010efficient}；改在学习者自己的策略产生的状态上训练、由专家标注，
可将其降到 $O(\epsilon T)$ \citep[DAgger;][]{ross2011dagger}。同策略蒸馏就是
以教师为专家的这套做法：学生写回合，教师在学生到达的每个状态上指出“换作我会
怎么做”。

\paragraph{取哪个方向：反向 KL}
目标函数是反向 KL，对学生自己的回合 $\tau$ 中的所有状态求和：
\begin{equation}
  \mathcal{L}(\theta) = \E_{x}\,\E_{\tau\sim\pis}\Big[\sum_{t}
  \KL\big(\pis(\cdot\mid s_t)\,\big\|\,\pit(\cdot\mid s_t)\big)\Big].
  \label{eq:objective}
\end{equation}
由 KL 的链式法则，这个求和的期望等于两个模型在完整回合上的分布之间的 KL；
工具结果在两个模型下都是历史的同一个函数，相互抵消。反向 KL 适合把大模型
蒸馏进小模型，理由有三。（i）它是学生分布下的期望，可以用学生的样本和教师对
这些 token 的对数概率来估计：教师只需一次前向传播，无须生成。（ii）它是寻模
（mode seeking）的：学生把概率放在教师几乎不放的地方会受重罚，漏掉教师的某些
选项只受轻罚。无法表示大教师全部行为的学生，会专注于一种好的行为，而不是把
几种行为平均起来 \citep{gu2024minillm,lu2025onpolicy}。（iii）只有当学生在其
访问的状态上表现得与教师一致时，它才达到最小值零，所以除了模仿，没有别的办法
能降低它。与学习得到的奖励模型不同，它无法被钻空子；\citet{lu2025onpolicy}
称之为 unhackable。

\paragraph{梯度就是策略梯度}
固定一个状态 $s$，记 $p=\pis(\cdot\mid s)$、$q=\pit(\cdot\mid s)$，$a$ 取遍
词表。由于 $p(a)\nabla_\theta\log p(a)=\nabla_\theta p(a)$ 且
$\sum_a\nabla_\theta p(a)=\nabla_\theta 1=0$，
\begin{align}
  \nabla_\theta \KL(p\,\|\,q)
  &= \sum_a \nabla_\theta p(a)\big[\log p(a)-\log q(a)\big]
    + \sum_a \nabla_\theta p(a) \nonumber\\
  &= \E_{a\sim p}\Big[\big(\log p(a)-\log q(a)\big)\,\nabla_\theta\log p(a)\Big].
  \label{eq:gradient}
\end{align}
对学生实际采样到的 token $y_t$，定义
\begin{equation}
  k_t = \log\pis(y_t\mid s_t) - \log\pit(y_t\mid s_t).
  \label{eq:k}
\end{equation}
由式 \eqref{eq:gradient}，$k_t\,\nabla_\theta\log\pis(y_t\mid s_t)$ 是 $s_t$
处 KL 梯度的单样本无偏估计。整个方法就是这一句：

\begin{keybox}{核心等式}
在学生自己的样本上最小化反向 KL \eqref{eq:objective}，就是 REINFORCE
\citep{williams1992reinforce}：最小化替代损失
\begin{equation}
  \ell(\theta) = -\frac{1}{N}\sum_{t} A_t\,\log\pis(y_t\mid s_t),
  \qquad A_t = -k_t \ \ \text{（视为常数）},
  \label{eq:loss}
\end{equation}
求和遍历一个批次中的 $N$ 个学生 token。token $t$ 的奖励为
\begin{equation}
  r_t = \log\pit(y_t\mid s_t) - \log\pis(y_t\mid s_t).
  \label{eq:reward}
\end{equation}
\end{keybox}

这个奖励等于教师对学生所写内容的对数概率，加上学生自己的自信息（surprisal，
相当于熵奖励）。教师比学生更认可的 token 被推高，不如学生认可的被压低。对 token
取平均后，$k_t$ 也是每 token 反向 KL 的无偏估计，即图~\ref{fig:curves} 中的
教师 KL 曲线。

\begin{algobox}[label=alg:step]{一个训练步}
\begin{algorithmic}[1]
\Require 学生 $\pis$（LoRA）、冻结的教师 $\pit$、环境、提示 $x_1,\dots,x_P$
\For{每个提示 $x_p$}
  \State 用 $\pis$（vLLM）采样 $G$ 个回合，执行每个工具调用并追加其结果
  \State $R_{p,i}\gets$ 回合 $i$ 的任务奖励；\quad
    $\hat A_{p,i}\gets(R_{p,i}-\bar R_p)/\sigma_p$
    \Comment{$G=1$ 时为 0}
\EndFor
\State 教师对每个回合做一次前向传播：
  $k_t\gets\log\pis(y_t\mid s_t)-\log\pit(y_t\mid s_t)$
\State 对回合 $(p,i)$ 的每个学生 token $t$：$A_t\gets c_R\,\hat A_{p,i}-c_T\,k_t$
\State 对 $\ell(\theta)=-\frac{1}{N}\sum_t A_t\log\pis(y_t\mid s_t)$
  做一步 AdamW
  \Comment{$N$ 个学生 token}
\end{algorithmic}
\smallskip
{\small\color{wr-taupe}各阶段只在环境、组大小 $G$ 和 $(c_T, c_R)$ 上不同。}
\end{algobox}

\begin{notebox}{两处近似}
第一，$\mathcal{L}$ 还通过状态依赖于 $\theta$：每个 token 都会影响其后的状态，
式 \eqref{eq:objective} 的精确梯度要在 $A_t$ 上再加上未来项
$\sum_{t'>t}k_{t'}$。我们与 \citet{lu2025onpolicy} 一样去掉这些项（折扣因子
取零）：每个 token 只为自己的分歧负责，这降低了方差，而下一批数据会在更新后的
学生所访问的状态上继续训练。第二，用一个采样 token 代替式 \eqref{eq:gradient}
中对整个词表的求和。GKD \citep{agarwal2024gkd} 精确计算这个和，需要两个模型在
每个位置上的完整下一 token 分布（Qwen3.5 每个位置 248{,}320 个数）；采样估计
只需要教师给出一个数。
\end{notebox}

\paragraph{为什么稠密信号胜过稀疏信号}
对比用二值任务奖励 $R\in\{0,1\}$ 的强化学习。GRPO 为每个提示采样一组 $G$ 个
回合，并给回合 $i$ 的每个 token 相同的优势 $\hat A_i=(R_i-\bar R)/\sigma_R$
（在组内归一化）。由此得出两点。一个回合的奖励至多携带一比特信息，还要分摊到
几百个 token 上：它只说明回合失败了，不说明错在哪里。而且，若一个提示的 $G$ 个
回合奖励全部相同，组内所有优势都为零。若学生以概率 $p$ 解出某个提示、且各样本
相互独立，则一组带有信号的概率为
\begin{equation}
  P_{\text{signal}}(p) = 1 - p^G - (1-p)^G ,
  \label{eq:signal}
\end{equation}
当 $G=4$ 时，它在 $p=\tfrac12$ 处为 0.875，在 $p=0.955$ 处为 0.17，$p$ 趋于 0
或 1 时趋于零。教师项则为每个 token 提供一个数，每个回合 $O(T)$ 比特而非
$O(1)$ \citep{lu2025onpolicy}；只要学生与教师有分歧它就不为零，无论回合成败。
它也很便宜：教师用一次前向传播（各位置并行）为整个回合打分，从不解码。

\paragraph{结合两种信号}
一般地，回合 $i$ 中 token $t$ 的优势为
\begin{equation}
  A_t = c_R\,\hat A_i - c_T\,k_t ,
  \label{eq:advantage}
\end{equation}
纯蒸馏时 $c_R=0$。教师项至多让学生与教师一样好；只有任务奖励能让学生超越教师
或纠正教师的错误，而且只能通过式 \eqref{eq:signal} 留下的那些有信号的组。
算法~\ref{alg:step} 是每个阶段的一个训练步。

\begin{notebox}{前提条件}
$k_t$ 比较的是两个模型对同一个 token id 的概率，所以学生和教师必须使用相同的
分词器：Qwen3（151{,}936 个 token）与 Qwen3.5（248{,}320 个）不能混用，这就是
每个系列只在自身内部蒸馏的原因。提示和工具结果进入 $s_t$，但不计入损失。此外，
教师必须评判它自己永远不会到达的状态，例如学生犯错之后的状态。
\end{notebox}

\section{训练设置}
\label{sec:setup}

\paragraph{模型与硬件}
表~\ref{tab:models} 列出两组教师–学生组合。Qwen3.5 \citep{qwen2026qwen35}
以三比一的比例交替使用 Gated DeltaNet 层 \citep{yang2025gateddeltanet}（一种
线性注意力）和门控注意力层。两次训练都使用非思考模式；在学生语言模型的每个
线性层上加秩为 32、$\alpha=32$ 的 LoRA \citep{hu2022lora}（Qwen3.5 的检查点
还带有视觉塔，这些纯文本任务用不到它）；都只用一块 Colab A100-SXM4-80GB。
训练器是 TRL 的 \texttt{GRPOTrainer} \citep{vonwerra2020trl}，用 vLLM 生成
\citep{kwon2023vllm}。我们的子类在每批采样之后计算 $k_t$，并用式
\eqref{eq:advantage} 替换 GRPO 的优势；采样、工具执行、掩码和损失都沿用 TRL。

\begin{table}[t]
\centering\small
\caption{两组教师–学生组合，其余一切相同。}
\label{tab:models}
\begin{tabular}{@{}lll@{}}
\toprule
\headrow & Qwen3 训练 & Qwen3.5 训练 \\
\midrule
教师（冻结，bf16） & Qwen3-8B & Qwen3.5-9B（仅语言模型） \\
学生 & Qwen3-1.7B & Qwen3.5-2B \\
学生层数 & 28 层，全为注意力 & 18 层 Gated DeltaNet + 6 层门控注意力 \\
词表大小 & 151{,}936 & 248{,}320 \\
工具调用格式 & JSON & XML \\
可训练 LoRA 参数 & 34.9\,M & 33.6\,M \\
\bottomrule
\end{tabular}
\end{table}

\paragraph{课程}
表~\ref{tab:stages} 列出四个阶段。阶段 1 采样 GSM8K 题目
\citep{cobbe2021gsm8k}，参考答案从不进入损失。阶段 2 的应用题数字足够大（比如
五位数乘四位数），用文字心算常常出错。阶段 3 和 4 在每个任务上都提供三种工具：
计算器、在随机生成的公司数据库上执行只读 SQL，以及在随机生成的商店文档上执行
JMESPath 查询。这些任务需要一种工具（45\%）、组合两种工具（40\%），或不需要
工具（15\%，应直接作答）。阶段 2--4 的任务由程序生成，训练与评估使用不同的
随机种子。每个阶段都在上一阶段的 LoRA 适配器上继续训练。

\begin{table}[t]
\centering\small
\caption{课程设置，两次训练完全相同。Token 上限：提示之后整个回合的长度上限，
含工具结果。批次：每步的回合数，即提示数 $\times$ 组大小 $G$。}
\label{tab:stages}
\begin{tabular}{@{}lllcrrrcc@{}}
\toprule
\headrow 阶段 & 任务 & 工具 & 轮数 & Token 上限 & 步数 & 批次 & 学习率
  & $(c_T, c_R)$ \\
\midrule
1 & GSM8K & 无 & 1 & 1{,}024 & 200 & $32\times1$ & $10^{-4}$ & (1, 0) \\
2 & 大数算术 & 计算器 & 4 & 2{,}048 & 150 & $32\times1$
  & $10^{-4}$ & (1, 0) \\
3 & 一种、两种或不用工具 & 全部三种 & 6 & 3{,}072 & 200 & $32\times1$
  & $10^{-4}$ & (1, 0) \\
4 & 同阶段 3 & 全部三种 & 6 & 3{,}072 & 150 & $16\times4$
  & $5\cdot10^{-5}$ & (1, 1) \\
\bottomrule
\end{tabular}
\end{table}

\paragraph{优化与打分}
\begin{itemize}[leftmargin=*,itemsep=1pt,topsep=2pt]
\item 每个优化步都使用由当前权重新采样的一批数据，因此 PPO 比率为 1，裁剪
  （下限 $\epsilon=0.2$，上限 0.28）从不起作用，更新就是式 \eqref{eq:loss} 的
  梯度，在批次内所有学生 token 上取平均（TRL 的 DAPO 损失；\citealp{yu2025dapo}）。
\item AdamW，5 步预热后学习率保持不变，梯度范数裁剪到 1，无权重衰减；每个微批
  8 个回合，使用梯度检查点。$10^{-4}$ 约为全参数微调学习率的十倍，这是 LoRA 的
  推荐做法 \citep{schulman2025lora}。
\item 采样：vLLM 与训练共用 GPU（占 30\% 显存），温度 1，top-$p$ 为 1。
\item 打分：教师对每个回合做一次 bf16 前向传播，经由融合的输出层直接返回采样
  token 的对数概率，不生成完整的 logits。$k_t$ 中学生的对数概率由训练器重新
  计算，而不取自 vLLM。
\item vLLM 的算子给出的对数概率与训练器略有差异，因此 TRL 给每个回合乘以重要性
  权重 $\exp\sum_t(\log\pis-\log\pi_{\text{vLLM}})$，并丢弃权重超过 3 的回合；
  各阶段的平均权重为 0.95--1.00。
\item 奖励（阶段 4 的训练与全部评估）：若最终回答中最后一个 \verb|\boxed{}| 与
  参考答案相等（数值误差 $10^{-6}$ 以内）则为 1，否则为 0；若回合用完 token
  上限，或到达轮数上限时仍在调用工具，也记为 0。
\item 每 25 步在 256 个留出任务上评估（阶段 1 用 GSM8K 测试集题目），每个任务
  以温度 1 采样一次，因为 TRL 没有贪心解码的评估。
\end{itemize}

\section{结果}
\label{sec:results}

表~\ref{tab:results} 给出每个阶段第一次和最后一次评估的结果；
图~\ref{fig:curves} 给出全部评估。

\begin{table}[t]
\centering\small
\caption{每个阶段开始和结束时的评估（256 个留出任务，温度 1），以及含评估在内的
实际耗时。}
\label{tab:results}
\setlength{\tabcolsep}{4.5pt}
\begin{tabular}{@{}lcccccc@{}}
\toprule
\headrow & \multicolumn{3}{c}{Qwen3-8B $\to$ Qwen3-1.7B}
 & \multicolumn{3}{c}{Qwen3.5-9B $\to$ Qwen3.5-2B} \\
\headrow 阶段 & 教师 KL & 奖励 & 耗时 & 教师 KL & 奖励 & 耗时 \\
\midrule
1. GSM8K & 0.229 $\to$ 0.136 & 0.809 $\to$ 0.836 & 45 分钟
  & 0.218 $\to$ 0.168 & 0.656 $\to$ 0.781 & 85 分钟 \\
2. 计算器 & 0.328 $\to$ 0.057 & 0.473 $\to$ 1.000 & 18 分钟
  & 0.108 $\to$ 0.047 & 0.645 $\to$ 0.902 & 103 分钟 \\
3. 三种工具 & 0.378 $\to$ 0.068 & 0.660 $\to$ 0.746 & 52 分钟
  & 0.163 $\to$ 0.034 & 0.754 $\to$ 0.957 & 109 分钟 \\
4. + 任务奖励 & 0.065 $\to$ 0.067 & 0.754 $\to$ 0.758 & 55 分钟
  & 0.034 $\to$ 0.025 & 0.945 $\to$ 0.953 & 130 分钟 \\
\midrule
合计 & & & 2.8 小时 & & & 7.1 小时 \\
\bottomrule
\end{tabular}
\end{table}

\begin{figure}[t]
\centering
% Evaluation curves of both runs, from data/*.dat (written by make_data.py).
% Shared by both reports: each defines the label macros \Lbl... before \input.
\pgfplotsset{
  run3/.style={color=wr-wine, mark=*, mark size=1.1pt, thick},
  run35/.style={color=wr-orange, mark=square*, mark size=1.1pt, thick},
}
% Stage numbers centred over each stage (x = 100, 275, 450, 625 of 700).
\newcommand{\stagenumbers}{%
  \node[font=\scriptsize, text=wr-taupe, anchor=north] at (rel axis cs:0.1429,1) {1};
  \node[font=\scriptsize, text=wr-taupe, anchor=north] at (rel axis cs:0.3929,1) {2};
  \node[font=\scriptsize, text=wr-taupe, anchor=north] at (rel axis cs:0.6429,1) {3};
  \node[font=\scriptsize, text=wr-taupe, anchor=north] at (rel axis cs:0.8929,1) {4};}
\begin{tikzpicture}
\begin{groupplot}[
  group style={group size=2 by 2, horizontal sep=1.5cm, vertical sep=1.45cm},
  width=0.5\textwidth, height=4.6cm,
  xmin=0, xmax=700, xtick={0,200,350,550,700},
  xmajorgrids,
  unbounded coords=jump,
  tick label style={font=\scriptsize},
  title style={font=\small\sffamily, text=wr-brown, yshift=-1ex},
  label style={font=\small},
]
\nextgroupplot[title={(a) \LblKL}, ymin=0, ymax=0.44,
  legend to name=curveslegend, legend columns=2,
  legend style={font=\small, draw=none,
    /tikz/every even column/.append style={column sep=1.5em}}]
\addplot[run3] table[x=step, y=kl] {data/qwen3.dat};
\addplot[run35] table[x=step, y=kl] {data/qwen35.dat};
\legend{\LblQwenThree, \LblQwenThreeFive}
\stagenumbers

\nextgroupplot[title={(b) \LblReward}, ymin=0.4, ymax=1.12]
\addplot[run3] table[x=step, y=reward] {data/qwen3.dat};
\addplot[run35] table[x=step, y=reward] {data/qwen35.dat};
\stagenumbers

\nextgroupplot[title={(c) \LblLength}, ymode=log, ymin=40, ymax=2000,
  xlabel={\LblStep}]
\addplot[run3] table[x=step, y=length] {data/qwen3.dat};
\addplot[run35] table[x=step, y=length] {data/qwen35.dat};

\nextgroupplot[title={(d) \LblEntropy}, ymode=log, ymin=5e-4, ymax=1,
  xlabel={\LblStep}]
\addplot[run3] table[x=step, y=entropy] {data/qwen3.dat};
\addplot[run35] table[x=step, y=entropy] {data/qwen35.dat};
\end{groupplot}
\end{tikzpicture}

\smallskip
\ref*{curveslegend}

\caption{四个阶段首尾相接训练期间的评估。虚线分隔各阶段；每个阶段开头的评估
用的是上一阶段的权重，但在新任务上进行。(a)~学生 token 上 $k_t$ 的均值。
(b)~256 个留出任务的解决比例。(c)~提示之后回合的平均长度（对数坐标）。
(d)~学生下一 token 分布的平均熵（对数坐标）。}
\label{fig:curves}
\end{figure}

\paragraph{蒸馏的效果与推导一致}
每个蒸馏阶段的教师 KL 都下降了 23\%--83\%，且大部分降幅出现在前 25--50 步
（Qwen3 阶段 2：到第 25 步已从 0.328 降至 0.049）。

\paragraph{不用奖励，任务成功率也随之上升}
阶段 1--3 从未见过奖励，但每个阶段的评估奖励都在上升：GSM8K 上分别提高 0.03 和
0.13，计算器任务上提高 0.53 和 0.26，三工具任务上提高 0.09 和 0.20（依次为
Qwen3 和 Qwen3.5）。学生学会了教师使用工具的习惯。在阶段 2 的第一批数据中，
Qwen3.5 学生在 32 个回合里只调用了 4 次计算器，其余算术都用文字完成；在最后
50 步中，它平均每个回合调用 1.55 次（Qwen3：第一批 32 个回合调用 15 次，最后
每回合 1.16 次）。在阶段 3，实际执行的工具调用中返回错误的比例从 21\% 降到
4\%（前五步对比最后五步；Qwen3：从 14\% 降到 7\%）。

\paragraph{拖累 Qwen3.5 的是长度，而不是准确率}
在 GSM8K 上，Qwen3.5 学生有 10\%--18\% 的回答超过 1{,}024 token 的上限而得零分。
这些并非死循环：遇到一个别扭的中间结果（比如算出的人数带小数）后，模型开始
反复自我怀疑。在写完的回答中，阶段 1 结束时有 91\% 是对的（Qwen3：84\%）。在
阶段 2 的最后 50 步中，它约 70\% 的失败训练回合用完了全部四轮，都在用更多的
计算器调用反复核对结果。

\begin{findbox}{任务奖励没有带来可测量的提升}
阶段 4 的评估奖励从 0.754 变为 0.758，从 0.945 变为 0.953。成功率约 0.95 时，
256 个任务的估计标准误为 0.014，而同一组 Qwen3.5 权重在阶段 3 结束时得分
0.957。式 \eqref{eq:signal} 解释了原因：只有 1.5\%（Qwen3）和 10.9\%
（Qwen3.5）的组存在奖励方差。若一个提示的四个结果是按平均成功率（0.80 和
0.955）独立抽取的，式 \eqref{eq:signal} 预测的比例应为 59\% 和 17\%。Qwen3
学生的结果远比这确定：它一旦做错某个提示，四个样本就全错。它的熵在阶段 2 崩塌（从
0.14 降到 0.002 nats/token），此后一直低于 0.05；Qwen3.5 学生从阶段 2 起保持在
0.13--0.16。这正是反向 KL 所预言的：在其最小值处，学生的下一 token 分布（连同
熵）等于教师的，所以自信的教师教出自信的学生（我们没有测量教师的熵），而自信的
学生让 GRPO 无从比较。Qwen3 训练的阶段 5（只用 GRPO，$G=8$，没有教师）证实了
这一点：第一批之后每个组的方差都为零，损失和梯度恰好为零，我们在大约第 20 步
停止了训练。
\end{findbox}

\paragraph{成本}
Qwen3.5 的训练用了 7.1 小时，Qwen3 用了 2.8 小时（约 50 对 23 个 Colab 计算
单元）。两次训练中每个 token 的成本相同：在阶段 4，按平均回合长度折算，每步
每 token 分别耗时 0.17 秒和 0.16 秒。但 Qwen3.5 学生写的回合长 2--5 倍，每次
调用工具前后都要用一句话解释（阶段 2：371 对 65 个 token；阶段 3：240 对 113）。

\section{要点}

\begin{tcbraster}[raster columns=2, raster equal height,
    raster column skip=4mm, raster row skip=4mm]
\begin{card}{wr-terracotta}{1\enspace 蒸馏是主力}
只用提示，没有标注也没有奖励，就能教会 1.7B 或 2B 的学生使用工具：每个阶段在
一块 A100 上不到两小时，大部分提升出现在前 50 步。
\end{card}
\begin{card}{wr-ochre}{2\enspace 强化学习需要能力边缘的任务}
GRPO 只能从各样本结果不一的提示中学习。蒸馏消除了随机性，而我们的任务大多已被
解决；更难的任务、更大的组或更高的温度，才是可用的杠杆。
\end{card}
\begin{card}{wr-wine}{3\enspace 预算要按模型定}
配置相同，但 Qwen3.5 写的回合长 2--5 倍：对 Qwen3 从未起作用的上限让它损失了
准确率，训练也用了 2.5 倍的 GPU 时间。
\end{card}
\begin{card}{wr-sage}{4\enspace 分词器决定谁能当教师}
教师必须与学生逐个 token 共享同一个词表，所以每个系列只能在自身内部蒸馏。
\end{card}
\end{tcbraster}

\Needspace{7\baselineskip}
\paragraph{复现}
阶段 1--4 用以下命令运行：
\begin{codebox}
BACKEND=trl EXTRA=-{}-use\_vllm scripts/run\_stages.sh a100 1 4
\end{codebox}
\noindent 训练器是 \path{opd/trl_train.py}，配置在 \path{opd/configs/}。所有记录的指标、
采样回合和适配器都在 \path{results/trl/a100/}（Qwen3）和
\path{results/trl/a100-qwen3.5/}（Qwen3.5）中。

\begin{thebibliography}{99}
\setlength{\itemsep}{1pt}
\small

\bibitem[Agarwal et~al.(2024)]{agarwal2024gkd}
R.~Agarwal, N.~Vieillard, Y.~Zhou, P.~Stanczyk, S.~Ramos, M.~Geist, and
O.~Bachem.
\newblock On-policy distillation of language models: Learning from
self-generated mistakes.
\newblock In \emph{ICLR}, 2024.

\bibitem[Cobbe et~al.(2021)]{cobbe2021gsm8k}
K.~Cobbe, V.~Kosaraju, M.~Bavarian, M.~Chen, H.~Jun, L.~Kaiser, M.~Plappert,
J.~Tworek, J.~Hilton, R.~Nakano, C.~Hesse, and J.~Schulman.
\newblock Training verifiers to solve math word problems.
\newblock arXiv:2110.14168, 2021.

\bibitem[Gu et~al.(2024)]{gu2024minillm}
Y.~Gu, L.~Dong, F.~Wei, and M.~Huang.
\newblock {MiniLLM}: Knowledge distillation of large language models.
\newblock In \emph{ICLR}, 2024.

\bibitem[Hu et~al.(2022)]{hu2022lora}
E.~J. Hu, Y.~Shen, P.~Wallis, Z.~Allen-Zhu, Y.~Li, S.~Wang, L.~Wang, and
W.~Chen.
\newblock {LoRA}: Low-rank adaptation of large language models.
\newblock In \emph{ICLR}, 2022.

\bibitem[Kim and Rush(2016)]{kim2016seqkd}
Y.~Kim and A.~M. Rush.
\newblock Sequence-level knowledge distillation.
\newblock In \emph{EMNLP}, 2016.

\bibitem[Kwon et~al.(2023)]{kwon2023vllm}
W.~Kwon, Z.~Li, S.~Zhuang, Y.~Sheng, L.~Zheng, C.~H. Yu, J.~E. Gonzalez,
H.~Zhang, and I.~Stoica.
\newblock Efficient memory management for large language model serving with
{PagedAttention}.
\newblock In \emph{SOSP}, 2023.

\bibitem[Lu and Thinking Machines Lab(2025)]{lu2025onpolicy}
K.~Lu and Thinking Machines Lab.
\newblock On-policy distillation.
\newblock \emph{Thinking Machines Lab: Connectionism}, 2025.
\newblock doi:10.64434/tml.20251026.

\bibitem[Qwen Team(2026)]{qwen2026qwen35}
Qwen Team.
\newblock {Qwen3.5} model cards.
\newblock Hugging Face, 2026.

\bibitem[Ross and Bagnell(2010)]{ross2010efficient}
S.~Ross and J.~A. Bagnell.
\newblock Efficient reductions for imitation learning.
\newblock In \emph{AISTATS}, 2010.

\bibitem[Ross et~al.(2011)]{ross2011dagger}
S.~Ross, G.~J. Gordon, and J.~A. Bagnell.
\newblock A reduction of imitation learning and structured prediction to
no-regret online learning.
\newblock In \emph{AISTATS}, 2011.

\bibitem[Schulman and Thinking Machines Lab(2025)]{schulman2025lora}
J.~Schulman and Thinking Machines Lab.
\newblock {LoRA} without regret.
\newblock \emph{Thinking Machines Lab: Connectionism}, 2025.

\bibitem[Shao et~al.(2024)]{shao2024deepseekmath}
Z.~Shao, P.~Wang, Q.~Zhu, R.~Xu, J.~Song, X.~Bi, H.~Zhang, M.~Zhang,
Y.~K. Li, Y.~Wu, and D.~Guo.
\newblock {DeepSeekMath}: Pushing the limits of mathematical reasoning in open
language models.
\newblock arXiv:2402.03300, 2024.

\bibitem[von Werra et~al.(2020)]{vonwerra2020trl}
L.~von Werra, Y.~Belkada, L.~Tunstall, E.~Beeching, T.~Thrush, N.~Lambert,
S.~Huang, K.~Rasul, and Q.~Gallou{\'e}dec.
\newblock {TRL}: Transformer reinforcement learning.
\newblock \url{https://github.com/huggingface/trl}, 2020.

\bibitem[Williams(1992)]{williams1992reinforce}
R.~J. Williams.
\newblock Simple statistical gradient-following algorithms for connectionist
reinforcement learning.
\newblock \emph{Machine Learning}, 8:229--256, 1992.

\bibitem[Yang et~al.(2025a)]{yang2025qwen3}
A.~Yang et~al.
\newblock {Qwen3} technical report.
\newblock arXiv:2505.09388, 2025a.

\bibitem[Yang et~al.(2025b)]{yang2025gateddeltanet}
S.~Yang, J.~Kautz, and A.~Hatamizadeh.
\newblock Gated delta networks: Improving {Mamba2} with delta rule.
\newblock In \emph{ICLR}, 2025b.

\bibitem[Yu et~al.(2025)]{yu2025dapo}
Q.~Yu et~al.
\newblock {DAPO}: An open-source {LLM} reinforcement learning system at
scale.
\newblock arXiv:2503.14476, 2025.

\end{thebibliography}

\end{document}
